\chapter{2016 年试题（当年科目代码 817）}

\begin{tishi}
\textbf{原卷卷面说明：}「试题附在答卷内交回，答案必须写在答题纸上，写在试题纸上无效。」
原卷封面所印科目代码 OCR 呈「817 工程流体力学」，本册按本考试科目通行代码 818 编排，请以原卷为准。
原卷共 5 页；下文各题〔图示〕均为依据 OCR 文字所作的解题用\textbf{文字描述}，
\textbf{不等于原卷图形}，\textbf{原卷图略处请对照原卷}，影印件见本册附录。
\end{tishi}

\section{一、选择题（共 30 分，每小题 3 分）}

\begin{ti}
  \item 液体的黏性主要来自于液体的（\quad）。\hfill\xstarfull{5}\yiju{E4 2026 计算 1（牛顿流体斜板静压强）；E5 题库选择 2（汽油＝牛顿流体）}\nandu{易}
  \fourchoices{分子热运动}{分子间的内聚力}{易变形性}{抗拒变形的能力}

  \item 相对压强的起算基准是（\quad）。\hfill\xstarfull{4}\yiju{E4 2018 选择 3（U 形水银测压计）；E4 2015 填空（U 形压差计）}\nandu{易}
  \fourchoices{绝对真空}{1 个标准大气压}{当地大气压}{液面压强}

  \item 液体在重力场中做加速直线运动时，其自由面与（\quad）处处正交。\hfill\xstarfull{3}\yiju{E6 教材 2.4；E5 题库}\nandu{中}
  \fourchoices{重力}{惯性力}{压力}{重力与惯性力的合力}

  \item 均匀流的（\quad）。\hfill\xstarfull{3}\yiju{E1 slide16 简答（均匀流的特点）}\nandu{易}
  \fourchoices{当地加速度为零}{迁移加速度为零}{向心加速度为零}{合加速度为零}

  \item 如图所示，在铅垂管中水自上而下流动，在相距 $l$ 的两断面间测得测压管水头差为 $h$，
  则两断面间的沿程水头损失 $h_f$ 为（\quad）。\hfill\xstarfull{5}\yiju{E4 2022 单选 3；E5 题库选择 10（圆管与方管沿程损失比 0.886）}\nandu{中}
  \fourchoices{$h_f=h$}{$h_f=h+l$}{$h_f=l-h$}{$h_f=h-l$}
  \tuyi{铅垂管路中水自上而下流动，在相距 $l$ 的两断面间接测压管，测得该两断面的测压管水头差为 $h$。原卷图略，请对照原卷；影印见本册附录。}
  \begin{tishi}选择题第 5 题的选项 D 原卷 OCR 为乱码（作 $h_f = =\ !$），
  按与 A、B、C 三项同一物理量关系补为 $h_f=h-l$。\textbf{原卷数值 OCR 不清。}\end{tishi}

  \item 圆管湍流过渡区的沿程阻力系数（\quad）。\hfill\xstarfull{3}\yiju{E6 教材 5.4；E2 slide5}\nandu{中}
  \fourchoices{与雷诺数 $Re$ 有关}{与管壁相对粗糙度 $\Delta/d$ 有关}{与 $Re$ 和 $\Delta/d$ 有关}{与 $Re$ 和管长 $l$ 有关}

  \item 如图所示，圆管层流流动过流断面上的切应力分布为（\quad）。\hfill\xstarfull{5}\yiju{E4 2022 单选 10（管径减半压强损失 16 倍）；E4 2015 填空（层流平均速度＝0.5 倍最大速度）}\nandu{中}
  \fourchoices{在过流断面上是常数}{管轴处是零，且与半径成正比}{管壁处是零，向管轴线性增大}{按抛物线分布}
  \tuyi{圆管层流流动，过流断面上的切应力分布；原卷 A、B、C、D 四个选项各配一幅切应力沿半径分布的示意图。
  原卷图略，请对照原卷；影印见本册附录。}

  \item 水流在管道直径、水温、沿程阻力系数都一定时，随着流量的增加，黏性底层的厚度就（\quad）。\hfill\xstarfull{3}\yiju{E6 教材 5.4；E2 slide5}\nandu{中}
  \fourchoices{增加}{减小}{不变}{不定}
  \begin{tishi}原卷此处写作「粘性底层」，本册统一为「黏性底层」，含义相同。\end{tishi}

  \item 图示两根完全相同的长管道，只是安装高度不同，两管道的流量关系为（\quad）。\hfill\xstarfull{5}\yiju{E5 题库计算 2（水箱垂直管道出流 $Q$ 与管长 $l$）；E1 slide33}\nandu{中}
  \fourchoices{$Q_1>Q_2$}{$Q_1<Q_2$}{$Q_1=Q_2$}{不定}
  \tuyi{两根完全相同的长管道（管长、管径、粗糙度均相同），仅安装高度不同，管内为长管出流。
  原卷图略，请对照原卷；影印见本册附录。}

  \item 孔口出流的流量系数、流速系数、收缩系数从小到大的正常排序是（\quad）。\hfill\xstarfull{5}\yiju{E1 slide33；E2 slide5；E4 多卷计算}\nandu{中}
  \fourchoices{流量系数、流速系数、收缩系数}{流量系数、收缩系数、流速系数}{流速系数、流量系数、收缩系数}{收缩系数、流速系数、流量系数}
\end{ti}

\section{二、填空题（共 30 分，每空 3 分）}

\begin{ti}
  \item 已知平面流动的速度场为 $u_x=x+2$，$u_y=-y+3$，则点 $(1,\,1)$ 的加速度
  $a_x=\underline{\hspace{2.5cm}}\ \mathrm{m/s^2}$，$a_y=\underline{\hspace{2.5cm}}\ \mathrm{m/s^2}$。\hfill\xstarfull{5}\yiju{E4 各年选择/填空（质点导数与加速度）；E1 slide33"公式较多，重点记忆"}\nandu{易}

  \item 用 U 形管测量 $A$ 处水压强，$h_1=0.15\,\mathrm{m}$，$h=0.3\,\mathrm{m}$，水银的重度 $\gamma_{Hg}=133280\,\mathrm{N/m^3}$，
  当时当地大气压强 $p_a=98000\,\mathrm{N/m^2}$，则 $A$ 处绝对压强等于 \underline{\hspace{2.5cm}} $\mathrm{N/m^2}$。\hfill\xstarfull{4}\yiju{E4 2018 选择 3（U 形水银测压计）；E4 2015 填空（U 形压差计）}\nandu{中}
  \tuyi{U 形水银测压计接在 $A$ 处（容器或管道）测压，测压计中水柱高度 $h_1$ 与水银面高差 $h$ 的具体标注位置见原卷图。
  原卷图略，请对照原卷；影印见本册附录。}
  \begin{tishi}水银重度原卷 OCR 作 $\mathrm{N/m^2}$，按重度的量纲改正为 $\mathrm{N/m^3}$；数值 $133280$ 未改动。\end{tishi}

  \item 长度 $l$、速度 $v$、重力加速度 $g$ 的无量纲组合是 \underline{\hspace{3cm}}；
  压强 $p$、密度 $\rho$、长度 $l$、流量 $Q$ 的无量纲组合是 \underline{\hspace{3cm}}。\hfill\xstarfull{4}\yiju{E1 slide33（量纲转化）；E2 slide5}\nandu{中}

  \item 在设计管道阀门实验时，首先要考虑满足的相似准则是 \underline{\hspace{2.5cm}} 规律。\hfill\xstarfull{4}\yiju{E4 2014 选择 10；E4 2016 填空 4；E4 2019 判断 2；E4 2022 单选 9；E5 题库选择 12}\nandu{中}
  \begin{tishi}原卷本空的断句 OCR 不清（作「……相似准则是（\quad）规律」），此处按 OCR 行序保留「规律」二字；
  请对照原卷核对空的位置与数量。\end{tishi}

  \item 在圆管流动中，层流的断面流速分布符合 \underline{\hspace{2.5cm}}。\hfill\xstarfull{5}\yiju{E4 2022 单选 10（管径减半压强损失 16 倍）；E4 2015 填空（层流平均速度＝0.5 倍最大速度）}\nandu{易}

  \item 密度 $\rho=0.85\,\mathrm{g/cm^3}$、运动黏度 $\nu=0.18\,\mathrm{cm^2/s}$ 的油在管径为 $100\,\mathrm{mm}$ 的管中
  以断面平均速度 $v=6.35\,\mathrm{cm/s}$ 作层流运动，管中心处的最大流速是 \underline{\hspace{2.5cm}} $\mathrm{cm/s}$，
  在离管中心 $r=20\,\mathrm{mm}$ 处的流速是 \underline{\hspace{2.5cm}} $\mathrm{cm/s}$，
  管壁处的切应力为 \underline{\hspace{2.5cm}} $\mathrm{N/m^2}$。\hfill\xstarfull{5}\yiju{E4 2022 单选 10（管径减半压强损失 16 倍）；E4 2015 填空（层流平均速度＝0.5 倍最大速度）}\nandu{中}
  \begin{tishi}原卷本空的分句（「……以断面平均速度 $v=$」「$=6.35\,\mathrm{cm/s}$ 作层流运动」「管中心处的最大流速是」）
  OCR 行序错乱，此处按物理含义归位，\textbf{数值未改动}。\end{tishi}
\end{ti}

\section{三、计算题（共 90 分）（1、2、3、4、5、6、7 题必做，8、9、10、11 题任选做两道）}

\begin{ti}
  \item 如图，一贮水容器，容器上有两个半球形的盖。设 $d=0.5\,\mathrm{m}$，$h=2.0\,\mathrm{m}$，$H=2.5\,\mathrm{m}$，
  求作用在两个球盖上的液体总压力。（12 分）\hfill\xstarfull{5}\yiju{E5 题库计算 1（弧形闸门 $F_x,F_z$ 与力矩）；E1 slide33}\nandu{中}
  \tuyi{一贮水容器上装有两个半球形盖，图中标注了球盖直径 $d$ 以及尺寸 $h$、$H$ 的基准位置。
  原卷图略，请对照原卷；影印见本册附录。}

  \item 如图，由水池通过等直径虹吸管输水。$A$ 点为虹吸管进口处，$H_A=0$；$B$ 点为虹吸管中与水池液面齐高的部位，
  $H_B=6\,\mathrm{m}$；$C$ 点为虹吸管中的最高点，$H_C=7\,\mathrm{m}$；$D$ 点为虹吸管的出口处，$H_D=4\,\mathrm{m}$。
  若不计流动中的能量损失，求虹吸管的断面平均流速和 $A$、$B$、$C$ 各断面上的绝对压强。（8 分）\hfill\xstarfull{5}\yiju{E4 2022 单选 6（计算点选取）；E4 2026 计算 2（文丘里流量证明）}\nandu{中}
  \tuyi{水池中插入一根等直径虹吸管：$A$ 为管子进口断面，$B$ 为管中与水池液面齐高处的断面，
  $C$ 为管子最高点处的断面，$D$ 为出口断面；图中给出各断面高程 $H_A=0$、$H_B=6\,\mathrm{m}$、$H_C=7\,\mathrm{m}$、$H_D=4\,\mathrm{m}$。
  原卷图略，请对照原卷；影印见本册附录。}
  \begin{tishi}原卷各断面高程符号下标 OCR 均作 $p$，此处按 $A$、$B$、$C$、$D$ 断面顺序改正为 $H_A$、$H_B$、$H_C$、$H_D$；数值未改动。\end{tishi}

  \item 用长 $l=800\,\mathrm{m}$、内径 $d=50\,\mathrm{mm}$ 的水平管道输油，若输油量要等于 $135\,\mathrm{L/min}$，
  求用以输油的油泵扬程为多大？（设油的密度 $\rho=920\,\mathrm{kg/m^3}$，动力黏度 $\mu=0.056\,\mathrm{Pa\cdot s}$）（8 分）\hfill\xstarfull{4}\yiju{E5 题库计算 4（水泵工况点与节流调节）；E6 教材 4.2}\nandu{中}
  \begin{tishi}原卷「用长＝800\,m」中长度符号 OCR 丢失，按题意补为 $l=800\,\mathrm{m}$；密度单位原卷 OCR 作 $\mathrm{kg/m^2}$，按密度改正为 $\mathrm{kg/m^3}$；数值均未改动。\end{tishi}

  \item 如图所示，水平方向的水射流以 $v_0=6\,\mathrm{m/s}$ 的速度冲击一斜置平板，射流与平板之间夹角 $\alpha=60^\circ$，
  射流过流断面面积 $A_0=0.01\,\mathrm{m^2}$，不计水流与平板之间的摩擦力，试求：
  （1）射流对平板的作用力 $F$；（2）流量 $Q_1$ 与 $Q_2$ 之比。（12 分）\hfill\xstarfull{5}\yiju{E4 2015 计算（水平弯管求 $F_x,F_y$）；E4 2026 计算 3（教材例题 4.12）；E1 slide16}\nandu{易}
  \tuyi{水平水射流以速度 $v_0$ 冲击一块斜置平板，射流与平板平面成 $\alpha$ 夹角，射流在平板上分为两股（流量分别记为 $Q_1$、$Q_2$），
  射流来流断面面积为 $A_0$。原卷图略，请对照原卷；影印见本册附录。}
  \begin{tishi}第（2）问原卷 OCR 作「流量 $Q$ 与 $Q_2$ 之比」，按射流冲击斜平板后分成两股的常规题意补为「流量 $Q_1$ 与 $Q_2$ 之比」，题设数值未改动。\end{tishi}

  \item 已知平面流动的速度分布为 $u_x=x^2-2x-4y$，$u_y=-2xy+2y$。试确定该流动：
  （1）是否满足连续性方程？（2）是否有旋？（3）如存在速度势函数 $\varphi$ 和流函数 $\psi$，求出 $\varphi$ 和 $\psi$。（10 分）\hfill\xstarfull{5}\yiju{E4 2015 计算（流函数求速度与流线）；E4 2026 计算 4（证明势函数存在性）}\nandu{中}
  \begin{tishi}原卷「$u_x\mathrm{-}x^2-2x-4y$」首字符 OCR 不清，按其与 $u_y$ 满足连续性方程的物理自洽性判读为 $u_x=x^2-2x-4y$；数值未改动。\end{tishi}

  \item 水从直径 $d$、长 $L$ 的铅垂管路流入大气中，水箱中液面高度为 $h$，管路局部阻力可忽略，
  沿程阻力系数为 $\lambda$。（10 分）\hfill\xstarfull{5}\yiju{E1 slide16"计算题：伯努利方程"；E1 slide33 与 E2 slide4"第四章＝考试计算题重点"}\nandu{中}
  \begin{xiaowen}
    \item 求管路起始断面 $A$ 处压强。
    \item $h$ 等于多少时，可使 $A$ 点的压强等于大气压。
  \end{xiaowen}
  \tuyi{水箱中的水经一根铅垂管路流入大气，水箱液面高度为 $h$，管路直径 $d$、长 $L$，$A$ 为管路起始断面。
  原卷图略，请对照原卷；影印见本册附录。}
  \begin{tishi}沿程阻力系数原卷 OCR 作「入」，按流体力学习惯改正为 $\lambda$。\end{tishi}

  \item 已知直径为 $15\,\mathrm{cm}$ 的输油管，流量为 $0.18\,\mathrm{m^3/s}$。现用流量为 $0.018\,\mathrm{m^3/s}$ 的水作模型实验，
  测得 $5\,\mathrm{m}$ 长输水管两端的压强水头差为 \underline{\hspace{2.5cm}} $\mathrm{m}$，求管两端的压强差 $\Delta p$（用油柱高表示）。（10 分）\hfill\xstarfull{3}\yiju{E4 2014 计算 8（溢流坝模型相似）；E5 题库选择 15}\nandu{难}
  \tuyi{原型为输油管，直径 $d=15\,\mathrm{cm}$、流量 $0.18\,\mathrm{m^3/s}$；模型为水，流量 $0.018\,\mathrm{m^3/s}$，测量段长 $5\,\mathrm{m}$。
  原卷图略，请对照原卷；影印见本册附录。}
  \begin{tishi}本空原卷 OCR 严重残缺：模型压强水头差的\textbf{数值无法判读}，
  \textbf{原卷数值 OCR 不清}；压强差表达式 OCR 仅存 $\Delta p$ 与 $\rho g$ 等碎片，已按题意排为「用油柱高表示」。\end{tishi}

  \item 如图所示，平面上有一对等强度的点涡，其强度为 $\Gamma$（$\Gamma>0$）、方向相反，分别位于 $(0,\,h)$、$(0,\,-h)$ 两固定点处；
  同时平面上有一无穷远平行于 $x$ 轴的来流 $v_x$。试求合成速度在原点的值。（10 分）\hfill\xstarfull{4}\yiju{E4 2015 计算 8（求驻点与流线）；E2 slide5}\nandu{难}
  \tuyi{平面上对称于原点布置一对方向相反、强度同为 $\Gamma$ 的点涡，分别位于 $(0,\,h)$ 与 $(0,\,-h)$ 两个固定点；
  另有平行于 $x$ 轴、来自无穷远的均匀来流 $v_x$。原卷图略，请对照原卷；影印见本册附录。}
  \begin{tishi}原卷 OCR 作「等强度为 $r$（$I>0$）的点涡」，按点涡强度的常用符号改正为 $\Gamma$（$\Gamma>0$）；
  来流符号原卷 OCR 作 $v_x$，保留原样。数值未改动。\end{tishi}

  \item 试确定某梯形断面均匀流渠道的水力最佳断面尺寸。已知边坡系数 $m=1.5$，糙率 $n=0.025$，底坡 $i=0.002$，
  流量 $Q=3.0\,\mathrm{m^3/s}$。（10 分）\hfill\xstarmark{X1}\nandu{难}
  \begin{tishi}明渠均匀流不在本考试第 1--8 章范围内，属旧大纲 / 复试范围设题（考点码 X1）。
  流量单位原卷 OCR 作 $\mathrm{m^2/s}$，按流量改正为 $\mathrm{m^3/s}$；数值 $3.0$ 未改动。\end{tishi}

  \item 一长为 $50\,\mathrm{m}$、浸水面积为 $469\,\mathrm{m^2}$ 的船以 $15\,\mathrm{m/s}$ 的速度在静水中航行，
  试求该船的摩擦阻力以及克服此阻力所需的功率。（设水的运动黏度 $\nu=0.011\,\mathrm{cm^2/s}$，摩擦阻力可按同一长度的相当平板计算）（10 分）\hfill\xstarfull{2}\yiju{E4 2015 计算 9（平板摩擦阻力比 $F_1/F_2$）}\nandu{难}
  \begin{tishi}运动黏度原卷 OCR 作 $0.01l\,\mathrm{cm^2/s}$，按常见取值的量级改正为 $\nu=0.011\,\mathrm{cm^2/s}$；
  \textbf{原卷数值 OCR 不清，请以原卷为准。}\end{tishi}

  \item 空气在管道内作恒定等熵流动，已知进口状态参数：$t_1=62\,^{\circ}\mathrm{C}$，$p_1=650\,\mathrm{kPa}$，$A_1=0.001\,\mathrm{m^2}$；
  出口状态参数：$p_2=452\,\mathrm{kPa}$，$A_2=5.12\times10^{-4}\,\mathrm{m^2}$。试求空气的质量流量 $Q_m$。
  （空气的气体常数 $R=287\,\mathrm{J/(kg\cdot K)}$，等熵指数 $k=1.4$）（10 分）\hfill\xstarmark{X2}\nandu{难}
  \begin{tishi}原卷 $A_2$ 的指数 OCR 缺失（作 $5.12\times10^{\sim}$），按其与 $A_1$ 的面积比与等熵流动质量守恒的自洽性
  判读为 $10^{-4}$，有效数字 $5.12$ 未改动；\textbf{原卷数值 OCR 不清，请以原卷为准}。\end{tishi}
\end{ti}
